\section{Descent of the common contact value}
\label{des:section}

The common contact value constructed in Section~\ref{ext:section}
is a section on a proper family over $X\times X$. We must show that its
value depends only on the two points of $X$. Properness first gives
finitely many possible values over each pair. The two vector lifts
then control these values along contact lines in opposite factors.
These tests exclude branching; purity and simple connectedness then
give a single value. We formulate this descent independently of the
stable-map construction. Throughout this section put
\[
Z=X\times X,\qquad M=L\boxtimes L.
\]

\begin{proposition}[Contact-value descent on the square]
\label{des:descent}
Let $X$ be a smooth simply connected complex projective contact Fano
manifold with $\Pic(X)=\ZZ[L]$. Assume that a family of free contact
lines covers a dense open subset of $X$. Let $p:Y\to Z$ be dominant,
where $Y$ is a proper integral algebraic space over $\CC$. Suppose that
there are regular sections
\[
W_{12}\in H^0\bigl(Y,p^*(T_X\boxtimes L)\bigr),\qquad
W_{21}\in H^0\bigl(Y,p^*(L\boxtimes T_X)\bigr)
\]
with the same contact value
\begin{equation}\label{des:common-value}
\xi=(\theta\boxtimes 1)(W_{12})
    =(1\boxtimes\theta)(W_{21})\in H^0(Y,p^*M).
\end{equation}
Then there is a unique section $B\in H^0(Z,M)$ such that $p^*B=\xi$.
\end{proposition}

\begin{proof}
\noindent\textbf{The finite image and its normalization.}
The section $\xi$ defines a morphism
$a:Y\to\Tot(M)$ over $Z$. It is proper: its graph is closed because
$\Tot(M)$ is separated over $\CC$, and projection from
$Y\times\Tot(M)$ is proper. Let $S_0$ be its reduced closed image.
This is an integral scheme, and $a$ factors through a proper surjection
$Y\to S_0$, since $Y$ is reduced.

The scheme $S_0$ is proper over $\CC$. Indeed it is separated and of
finite type, and universal closedness descends along the surjection
$Y\to S_0$: after any base change, a closed subset of $S_0$ has
closed inverse image in $Y$ with the same image in the base. Since
$S_0$ is closed in the line-bundle total space, $q_0:S_0\to Z$ is
affine. It is also proper, hence finite. Dominance of $p$ makes $q_0$
surjective.

Let $\nu:S\to S_0$ be the normalization and put $q=q_0\nu$.
A scheme of finite type over $\CC$ is Nagata, so normalization is
finite \cite[Tags \stacksTag{035B} and \stacksTag{035S}]{Stacks}. Thus $q$ is a finite
dominant morphism of integral schemes, $S$ is normal, and $q$ is
generically \'etale by characteristic zero.
We require an open set in the \emph{base} on which normalization
causes no change. The morphism $\nu$ is an isomorphism over a dense
open $V\subset S_0$. The finite image $q_0(S_0\setminus V)$ is a
proper closed subset of $Z$, since finite maps preserve the dimension
of closed subsets. The image under $q$ of its non-\'etale locus is
likewise proper and closed. Deleting both images gives a dense open
\begin{equation}\label{des:base-open}
U\subset Z\quad\text{such that}\quad
S\times_ZU\simeq S_0\times_ZU\longrightarrow U
\text{ is finite \'etale}.
\end{equation}
No morphism from $Y$ to the normalization $S$ is required.

\medskip
\noindent\textbf{Every generic sheet splits over either type of line.}
Choose a smooth irreducible open family $H$ of parametrized free contact
lines with smooth dominant evaluation
\[
\operatorname{ev}:H\times\PP^1\longrightarrow X.
\]
This is possible by Corollary~\ref{lines:free-cover}: freeness gives
$H^1(\PP^1,f^*T_X(-t))=0$ for every $t$, so evaluation is smooth at
every pair $(f,t)$ on the chosen open of the Hom scheme. For each
$f\in H$, let $P_f$ be the positive part of $f^*T_X$. Thus
\begin{equation}\label{des:line-quotient}
f^*T_X/P_f\simeq\OO_{\PP^1}^{\oplus r_f},
\qquad \theta(P_f)=0.
\end{equation}
The second assertion is Lemma~\ref{lines:positive}.


Consider a map
\[
c:\PP^1\longrightarrow Z,\qquad c=(f,y),\quad f\in H,
\]
whose image meets $U$. Set $T=\PP^1\times_ZY$, and let
$D\subset T$ be any reduced irreducible component dominating
$\PP^1$, with projection $\pi_D:D\to\PP^1$. It is a proper
integral algebraic space over $\CC$, and
$H^0(D,\OO_D)=\CC$. For the latter assertion, a regular function
gives a morphism to $\PP^1$ whose proper irreducible image avoids
infinity, hence is a point.

Pull $W_{12}$ back to $D$ and project to
\[
\pi_D^*\bigl((f^*T_X/P_f)\otimes L_y\bigr).
\]
In a global trivialization this section is a constant vector, by
\eqref{des:line-quotient} and the preceding assertion about functions.
The contact map factors through this quotient. Consequently there
exists a section
\begin{equation}\label{des:component-section}
a_D\in H^0(\PP^1,c^*M),\qquad
\xi|_D=\pi_D^*a_D.
\end{equation}
For the other type $c=(y,f)$, the identical argument uses $W_{21}$
and proves the same assertion.

We spell out why this accounts for \emph{all} sheets. Write
$K=\CC(\PP^1)$ and $\eta=\Spec K$. The surjection $Y\to S_0$
remains surjective after base change, so
\[
T_\eta\longrightarrow (\PP^1\times_ZS_0)_\eta
\]
is surjective. There are finitely many irreducible components of $T$;
the components that do not dominate $\PP^1$ disappear over $\eta$.
Every point of $T_\eta$ therefore belongs to the generic fiber of one
of the components $D$ just considered. Trivialize $c^*M$ over $K$.
Then $(\PP^1\times_ZS_0)_\eta$ is a finite closed subscheme of
$\mathbb A^1_K$, and \eqref{des:component-section} puts the image of
$D_\eta$ at the $K$-rational point whose coordinate is $a_D(\eta)$.
Indeed the induced map from the coordinate algebra $K[\zeta]$ sends
$\zeta$ to the scalar $a_D(\eta)\in K$, so it factors through
$K[\zeta]/(\zeta-a_D(\eta))=K$. This does not require a
$K$-rational point of $D_\eta$.
Surjectivity shows that every point of this finite subscheme is
$K$-rational. By \eqref{des:base-open} the subscheme agrees with
$(\PP^1\times_ZS)_\eta$ and is finite \'etale, hence reduced.
Its algebra is therefore a product of copies of $K$.

In particular, for \emph{every} $c$ of either type that meets $U$,
\begin{equation}\label{des:all-sheets}
\begin{gathered}
C\subset(\PP^1\times_ZS)_{\mathrm{red}}\text{ an irreducible component},
\quad C\longrightarrow\PP^1\text{ dominant}\\
\Longrightarrow\qquad \deg(C^\nu/\PP^1)=1,
\end{gathered}
\end{equation}
where $C^\nu$ denotes the normalization of $C$.
This conclusion uses neither flatness of $T$ nor connectedness or
reducedness of its fibers. The embedding of $S_0$ in a line bundle is
what converts the componentwise contact values into rational sheets.

\medskip
\noindent\textbf{The local model at a ramification divisor.}
Suppose that $q$ is not unramified at the generic point of some
prime divisor $R\subset S$. Its image is a prime divisor
$\Delta\subset Z$, since $q$ is finite. Put $N=\dim Z=2\dim X$.
There is a dense smooth open $\Delta^0\subset\Delta$ such that every
point above $\Delta^0$ is smooth in $S$, and the reduced divisors
$R_i$ above $\Delta$ are smooth, disjoint, and \'etale over
$\Delta^0$. To obtain this open, delete from $\Delta$ the finite image
of $\operatorname{Sing}(S)\cap q^{-1}(\Delta)$ and those of the
singular loci of the $R_i$, their pairwise intersections, and the
non-\'etale loci of $R_i\to\Delta$, as well as the singular locus
of $\Delta$. These are proper closed subsets
of $\Delta$: normality makes $S$ regular at every divisorial generic
point, and the finite maps of reduced divisors are generically
separable in characteristic zero.

At a point $s\in R_i$ over $x\in\Delta^0$, choose local analytic
coordinates $(x_1,\ldots,x_N)$ on $Z$ with $\Delta=(x_1=0)$.
If $z_1$ locally defines $R_i$, then
$q^*x_1=u z_1^{e_i}$ for a unit $u$. Absorb an analytic $e_i$-th
root of $u$ into $z_1$ and take $q^*x_2,\ldots,q^*x_N$ as the
remaining coordinates; their restrictions are coordinates along
$R_i$ because its map to $\Delta^0$ is \'etale. The map becomes
\begin{equation}\label{des:branch-model}
(z_1,\ldots,z_N)\longmapsto
(z_1^{e_i},z_2,\ldots,z_N).
\end{equation}
At least one $e_i$ is greater than one. Otherwise this model would
make $q$ unramified at every generic point over $\Delta$, contrary
to the choice of $R$. Each point of $\Delta^0$ has a preimage on
such a ramified $R_i$, since $R_i\to\Delta$ is finite surjective.

\medskip
\noindent\textbf{A test line detects every possible divisor.}
Let $d=\dim X$. If $\Delta\to X$ under the second projection is
dominant, then for general $y\in X$ its intersection with
$X\times\{y\}$ is a nonempty effective Cartier divisor
$\Delta_y$ on $X$. It is Cartier by restricting a local equation of
$\Delta$ to a fiber not contained in $\Delta$, and its nonemptiness
and dimension $d-1$ follow from properness and the fiber dimension
Theorem. Write $\OO_X(\Delta_y)=L^{\otimes a}$. A nonzero effective
divisor has positive intersection with $L^{d-1}$, so $a>0$.
For every contact line $f$ not contained in $\Delta_y$,
\[
\deg f^*\OO_X(\Delta_y)=a\deg f^*L=a>0.
\]
Thus a general test line of the first type $(f,y)$ meets $\Delta$.

If $\Delta$ does not dominate the second factor, its closed image
$D\subset X$ has dimension at least $d-1$, because
$\dim\Delta=2d-1$ and its fibers have dimension at most $d$.
Hence $D$ is a prime divisor, and the inclusion
$\Delta\subset X\times D$ between irreducible varieties of the
same dimension is equality. In this case $\OO_X(D)=L^{\otimes a}$
with $a>0$ by the same intersection argument, and a general line of
the second type $(y,f)$ meets $\Delta=X\times D$.

For completeness, the required general-position assertions concern
the whole chosen line family, whose evaluation need only be dominant.
For either type let $P=H\times X$, put $m=\dim P$, and denote the
universal evaluation by
\[
E:P\times\PP^1\longrightarrow Z.
\]
It is smooth and dominant. If $A=\Delta\setminus\Delta^0$, then
$A$ is closed of codimension at least two in $Z$ and
\[
\dim E^{-1}(A)\leq m-1.
\]
Projection to $P$ is proper, so the image of $E^{-1}(A)$ is a
proper closed subset of $P$. A general test line therefore avoids
$A$. The scheme $E^{-1}(\Delta^0)$ is smooth of pure dimension $m$.
Discard the closures of the images of its nondominating components,
and apply generic
smoothness to the projections of its dominating components to $P$.
For general parameters their fibers are smooth of dimension zero.
These fibers are the inverse images of $\Delta^0$ on the test line,
so all intersections are transverse. Finally, meeting $U$ and not
being contained in $\Delta$ are nonempty open conditions on $P$:
they are images under the open projection $P\times\PP^1\to P$ of
the corresponding nonempty open inverse images under $E$.
All these open conditions hold simultaneously because $P$ is
irreducible. For the type selected in the preceding two paragraphs,
positive intersection ensures at least one intersection with
$\Delta^0$.

\medskip
\noindent\textbf{Ramification contradicts the split sheets.}
Choose this test line $c$, a transverse intersection $c(t_0)\in
\Delta^0$, and a ramified point $s$ above it. In
\eqref{des:branch-model} let $e=e_i>1$ and take a local parameter
$t$ at $t_0$. Eliminating the tangential coordinates, the analytic
local ring of $\PP^1\times_ZS$ at $(t_0,s)$ is
\begin{equation}\label{des:pullback-germ}
\CC\{z_1,t\}/(z_1^e-a(t)),\qquad
a(0)=0,\quad a'(0)\ne0.
\end{equation}
It is isomorphic to $\CC\{z_1\}$ by the analytic inverse function
Theorem. Thus this pullback germ itself is smooth, reduced, and
irreducible, and $t$ has order $e$ on it. In particular no vertical
or embedded component passes through this germ, whatever happens
elsewhere in the finite pullback. The unique reduced irreducible
component through it dominates $\PP^1$. Its normalization is
unchanged near this smooth point, and the resulting finite map of
normal curves has ramification index $e$. The degree of that map is
at least $e$, by the degree formula for a fiber of a finite map of
smooth curves. This contradicts \eqref{des:all-sheets}. We have proved
that $q$ is unramified at every codimension-one point of $S$.

\medskip
\noindent\textbf{Purity and simple connectedness.}
We check the local hypotheses of purity at an arbitrary nongeneric
point $s\in S$, with $z=q(s)$. Both schemes are Noetherian,
$\OO_{S,s}$ is normal, $\OO_{Z,z}$ is regular, and $q$ is
quasi-finite at $s$ because it is finite. Moreover
\[
\dim\OO_{S,s}=\dim\OO_{Z,z}>0.
\]
Indeed, on affine charts the corresponding finite extension of
domains has going down because the base domain is integrally closed
\cite[Proposition 10.38.7, Tag \stacksTag{00H8}]{Stacks}.
Going down and incomparability of primes in an integral extension
give equality of the heights of corresponding primes, including
nonclosed points. For every generization $s'$ of $s$ with
$\dim\OO_{S,s'}=1$, the map $q$ is unramified at $s'$ by the
preceding paragraph. These are precisely the local hypotheses of
purity of the branch locus, which makes $q$ \'etale at $s$
\cite[Lemma 58.21.4, Tag \stacksTag{0BMB}]{Stacks}. The generic point
is already \'etale by separability. In particular flatness is a
conclusion here, not a premise about $q$ or its line pullbacks.

The analytic space $Z^{\mathrm{an}}=X^{\mathrm{an}}\times
X^{\mathrm{an}}$ is simply connected. The finite \'etale map $q$
analytifies to a finite covering, which is connected because $S$ is
integral. It has degree one. Consequently $S_0\to Z$ is finite
birational and hence an isomorphism, since $Z$ is normal
\cite[Lemma 29.55.8, Tag \stacksTag{0AB1}]{Stacks}. Its closed inclusion into
$\Tot(M)$ is the graph of a global section $B$ with $p^*B=\xi$.
Finally, pullback of sections of a line bundle along a dominant map
of integral spaces is injective. This proves uniqueness.
\end{proof}

\begin{corollary}[Symmetry and the full local family]
\label{des:family}
Apply Proposition~\ref{des:descent} to the space and tensors of
Proposition~\ref{ext:moving-pair}. The resulting section
$B\in H^0(X\times X,L\boxtimes L)$ is symmetric under exchange of
the factors, using the canonical exchange of the two line factors.
For the full neighborhood $\mathcal H$ of $g$ in the space of
parametrized degree-two maps provided there, one has
\begin{equation}\label{des:family-identity}
B\bigl(h(z),h(w)\bigr)
=(z-w)^2\,\theta\bigl(dh(\partial_z)\bigr)
       \otimes\theta\bigl(dh(\partial_w)\bigr),
\qquad h\in\mathcal H.
\end{equation}
This identity holds for all pairs of source points, including the
diagonal; its displayed expression uses any common affine coordinate.
\end{corollary}

\begin{proof}
Let $\sigma:Z\to Z$ exchange the two factors, and let
$\iota:Y\to Y$ be the involution in
Proposition~\ref{ext:moving-pair}. Under the canonical line-bundle
identifications, $p\iota=\sigma p$ and $\iota^*\xi=\xi$. Hence
\[
p^*(\sigma^*B)=\iota^*(p^*B)=\iota^*\xi=\xi=p^*B.
\]
Dominance of $p$ gives $\sigma^*B=B$.

The equality $p^*B=\xi$ pulls back along the morphism
$\mathcal H\times\PP^1\times\PP^1\to Y$ supplied by
Proposition~\ref{ext:moving-pair}. Its smooth-source formula
\eqref{ext:hom-identity} gives \eqref{des:family-identity}.
Here $\mathcal H$ is a neighborhood in the full Hom space by
Lemma~\ref{sm:full-germ}, so no deformation directions at $g$ are lost.
Coordinate independence
and regularity at infinity follow from the invariant tensor
$(z-w)^2\partial_z\otimes\partial_w$; its double zero on the
diagonal also proves the asserted diagonal identity.
\end{proof}
